% =====================================================================
% Appendix: continuous and graded metrics.
%
% CREATED 2026-08-13, in response to reviewer feedback ("Methodological &
% Scope Clarifications"): the only acknowledgement of float-scored metrics was
% one clause in sections/limitations.tex, and the reviewer asked for a formal
% treatment. The body keeps a compact subsection
% (\S\ref{sec:limitations:continuous}); the algebra and the per-metric mapping
% live here, because the body is at its measured page target
% (docs/PAGE_MAP_2026-08-05.md, limitations 2 pages against a target of 1.5).
%
% STATUS RULE FOR THIS WHOLE FILE, AND IT IS THE POINT OF IT.
% Nothing here was run. Everything in it is either (a) MATHEMATICS -- the
% standard TOST planning argument applied to the mean of a bounded paired
% difference, which is a derivation and not an empirical claim -- or (b) a
% PROPOSAL for how a metric family would map onto that argument. No required-n
% in this paper transfers to any non-binary metric, and this appendix must never
% be written as though the extension had been validated. If a later revision
% measures a sigma_d for any metric, that measurement gets its own registration
% and its own SOURCE pointer; it does not get folded into this file's prose.
%
% NUMBERS. This file states NO empirical number of its own. The two it points at
% are consumed through macros or cross-references:
%   \AuditOutsideFrameworkClaim{}, \AuditPerItemTaskMatched{}, \AuditEligible{}
%     -> paper/audit_denominators.tex (GENERATED from the sealed
%        results/audit_verdicts_rev3.csv)
%   the excluded float-scored atlas cells
%     -> Appendix~\ref{app:atlas:exclusions}, which states the count and its
%        source; it is deliberately not restated here (prose rule 9).
%
% WORDING RULE, inherited from sections/audit.tex §\ref{sec:audit:r04} and the
% carry checklist §15: R04 is outside OUR REGISTERED BINARY PAIRED-OUTCOME
% CALCULATION. NEVER write "incompatible with a paired framework" -- CIDEr
% supports paired resampling perfectly well; it is the flip model that does not
% apply. This file must not state the audit's verdict for R04 in any form the
% audit section does not use.
% =====================================================================

\section{Continuous and graded metrics: the extension this paper does not run}
\label{app:continuous}

This appendix states what would be required to certify a metric that scores each
item on a scale instead of marking it right or wrong. It is a derivation and a
proposal. No part of it has been executed, no quantity in it has been measured,
and none of the required sample sizes reported in this paper applies to any
metric outside the binary per-item model.

\subsection{Why the binary outcome model comes first}
\label{app:continuous:why}

Three things follow from binary per-item correctness, and the certification
apparatus uses all three.

The first is an exact variance. Write $c_i^{\mathrm{base}}, c_i^{\mathrm{comp}}
\in \{0,1\}$ for the correctness of item $i$ under the baseline and the
compressed model, and $d_i = c_i^{\mathrm{comp}} - c_i^{\mathrm{base}}$ for the
per-item accuracy difference, so that $d_i \in \{-1,0,+1\}$. Because $d_i^2$ is
the indicator that the two models disagree on item $i$,
\begin{equation}
\mathrm{Var}(d_i) \;=\; \mathbb{E}[d_i^2] - \delta^2 \;=\; p_d - \delta^2 ,
\qquad \delta = \mathbb{E}[d_i],
\label{eq:vard-binary}
\end{equation}
and under the null of no true difference this is $p_d$ exactly, which is the
$\mathrm{sd}_{\mathrm{paired}}$ of Equation~\eqref{eq:sds}. Nothing is assumed
about the joint distribution of the two runs beyond the pairing itself: the
second moment of a variable supported on $\{-1,0,+1\}$ \emph{is} the probability
that it is non-zero. That identity is what makes the certification tables
countable. The variance term is a rate a third party obtains by comparing two
per-item dumps, so Table~\ref{tab:certification} is populated by measurement
(\S\ref{sec:atlas}) and not by an assumption about how models vary.

% SOURCE (non-numeric): the mined quantity is the per-item correctness field,
% docs/ATLAS_MINING_REGISTRATION_2026-07-15.md §§1, 3.4 as reproduced in
% Appendix~\ref{app:atlas:sources}; the excluded float-scored cells and their
% count are Appendix~\ref{app:atlas:exclusions}, whose SOURCE pointer is
% results/atlas_cells_summary_rev2.csv (reason column). The count is
% deliberately not restated here.
The second is that binary correctness is what the public record carries. The
atlas mines per-item dumps whose recorded quantity is a correctness field
(Appendix~\ref{app:atlas:sources}); the cells it met that carried no binary
correctness metric at all were excluded rather than scored, and that exclusion
is reported with its count in Appendix~\ref{app:atlas:exclusions}. A continuous
apparatus built on this record would have had nothing to be built from.

% SOURCE: \AuditNotAssessableOutsideFramework{} and
% \AuditOutsideFrameworkClaim{} are generated from
% results/audit_verdicts_rev3.csv via paper/audit_denominators.tex. This states
% the audit's own count and does not generalise it into a claim about which
% metrics the sources report.
The third is that binary correctness is the outcome type the audit could score.
Of the eligible claims, exactly \AuditNotAssessableOutsideFramework{} was set
aside because the metric its qualifying sentence speaks about carries no
per-item correctness state, and that claim, \AuditOutsideFrameworkClaim{}, is
the subject of \S\ref{app:continuous:r04} below.

\subsection{TOST for a bounded graded score}
\label{app:continuous:tost}

The generalisation of the planning calculation is immediate, and stating it
costs nothing, because TOST is defined on the mean of a paired difference and
never on its support. Let each item carry bounded scores $s_i^{\mathrm{base}},
s_i^{\mathrm{comp}} \in [a,b]$ under the same metric, put $d_i =
s_i^{\mathrm{comp}} - s_i^{\mathrm{base}}$, and let the reported statistic be
$\bar{d} = n^{-1} \sum_i d_i$ with $\sigma_d^2 = \mathrm{Var}(d_i)$. TOST at
one-sided $\alpha$ with power $1-\beta$, against an assumed true difference of
zero and at a margin $m$ expressed in the metric's own units, then requires
\begin{equation}
n_{\mathrm{req}} \;=\; \left\lceil \left(\frac{(z_{1-\alpha} + z_{1-\beta})\,
\sigma_d}{m}\right)^{\!2} \right\rceil ,
\label{eq:nreq-continuous}
\end{equation}
which is Equation~\eqref{eq:nreq} with the paired standard deviation replaced by
$\sigma_d$. The binary case is the special case $\sigma_d^2 = p_d$ of
Equation~\eqref{eq:vard-binary}. Two consequences are pure algebra and worth
stating because they bound what any adaptation can achieve: the requirement
scales as $(\sigma_d/m)^2$, so a graded metric needs more items than the
corresponding binary row of Table~\ref{tab:certification} exactly when
$\sigma_d^2 > p_d$ and fewer when it is smaller; and every qualification already
attached to Equation~\eqref{eq:nreq} carries over unchanged, including that this
is a planning size under an assumed true difference of zero and is therefore a
lower bound on what a non-zero true difference would demand
(\S\ref{sec:cert:method}).

What does not transfer is everything empirical.

\emph{The churn tables do not transfer.} $\sigma_d$ is a property of the metric,
of the model pair and of the benchmark family jointly, and it has to be measured
for each combination in the same way $p_d$ was. This paper measures it for none.
The atlas records accuracy-state churn and carries no per-item score dispersion
for any metric, so there is no path from the numbers in this paper to a
$\sigma_d$ for CIDEr, for ROUGE, or for a judge scale. A practitioner adapting
the method supplies that measurement; the method does not supply it to them.

\emph{The interpretation of churn does not transfer.} For a binary outcome,
churn is a count: the proportion of items whose correctness state changed, which
is simultaneously the variance term and a directly reportable description of
per-item behaviour. For a graded metric the corresponding quantity is a
dispersion of per-item score changes, and it is not a count of anything. In
particular the net-versus-gross decomposition that motivates this paper has to
be redefined before it means anything: a cell can show a net difference of zero
with arbitrarily large $\sigma_d$, but "how many items changed" is no longer the
question, since almost every item changes a little.

\emph{The margin does not transfer.} $m$ in
Equation~\eqref{eq:nreq-continuous} is in the metric's own units. A margin of two
CIDEr points and a margin of two accuracy points are unrelated quantities, and
neither can be read across to the other, so a certificate on a graded metric is
interpretable only if the margin is declared in that metric and defended in that
metric. The registered $2$\,pp margin of this paper is an accuracy margin and
means nothing elsewhere.

\subsection{How three metric families would map}
\label{app:continuous:families}

Each of the following is a proposal about structure, not a result.

\paragraph{pass@$k$.} At $k=1$ the per-item outcome is a single pass or fail, so
this is the binary model already, with the added wrinkle that sampling makes
correctness itself random and one draw is one draw. For $k > 1$ the per-item
quantity is an estimated success probability $\hat{\pi}_i \in [0,1]$, computed
from a finite number of completions. The per-item difference is then bounded and
graded, so Equation~\eqref{eq:nreq-continuous} applies in form, but the
estimation error enters the variance: writing $\pi_i$ for the quantity being
estimated, the law of total variance gives
\begin{equation}
\sigma_d^2 \;=\; \mathrm{Var}\!\left(\pi_i^{\mathrm{comp}} -
\pi_i^{\mathrm{base}}\right) \;+\;
\mathbb{E}\!\left[\mathrm{Var}\!\left(\hat{\pi}_i^{\mathrm{comp}} -
\hat{\pi}_i^{\mathrm{base}} \,\middle|\, \pi_i^{\mathrm{comp}},
\pi_i^{\mathrm{base}}\right)\right],
\label{eq:passk-variance}
\end{equation}
whose second term shrinks as the number of completions per item grows. Sizing a
pass@$k$ evaluation therefore involves two budgets, items and completions per
item, and only the first is what a required-$n$ table counts. If the same
sampling seeds drive both sides, the two estimation errors are positively
correlated and the second term is smaller; that is a design choice to be
declared, not a property to be assumed.

\paragraph{BLEU, ROUGE and CIDEr.} These assign a bounded score to a generated
string, so at the level of a single item they are exactly the graded case above
and Equation~\eqref{eq:nreq-continuous} applies once $\sigma_d$ is measured. The
obstacle is aggregation rather than gradedness. Corpus-level BLEU is not the
mean of per-sentence BLEU: it is computed from n-gram match counts and reference
lengths summed over the whole corpus, with a brevity penalty formed from those
totals, so the reported statistic is not an average of per-item quantities and
there is no $d_i$ whose mean it is. Certifying it therefore requires either a
sentence-level variant, which is a different number from the corpus figure and
must be declared as the certified quantity, or a paired resampling scheme over
items that respects the aggregation, in which case the sample size is found by
simulation and Equation~\eqref{eq:nreq-continuous} does not give it. ROUGE and
CIDEr are ordinarily reported as means of per-item scores and do not have this
particular problem, which is why the aggregation question has to be asked per
metric rather than per family.

\paragraph{LLM-as-a-judge scales.} A judge produces a bounded ordinal score, so
the graded form applies in principle, with two additions. First, the score is
itself a model output, so the observed per-item value is the quantity of
interest plus judge error. Under the simple model that an observed score is
$s_i + e_i$ with judge errors of variance $\sigma_j^2$ and correlation $\rho$
between the two sides, the difference has variance
\begin{equation}
\sigma_d^2 \;=\; \sigma_{\mathrm{true}}^2 \;+\; 2\sigma_j^2 (1 - \rho),
\label{eq:judge-variance}
\end{equation}
so judge noise inflates the requirement unless the same judge scores both sides
of an item in a way that makes its errors correlated. $\sigma_j^2$ has to be
estimated from repeated judgements, which is a third budget, and assuming it
away is the failure mode this paper exists to complain about in another setting.
Second, an ordinal scale is not an interval scale, so a margin stated as a
fraction of a point on a five-point rubric has no agreed meaning; the margin
would have to be defended as a statement about the rubric before the arithmetic
is worth running. We take no position here on how that defence should go.

\subsection{\AuditOutsideFrameworkClaim{} and what scoring it would require}
\label{app:continuous:r04}

\AuditOutsideFrameworkClaim{} is the one audited claim that this limitation
touches directly, and the audit's treatment of it stands as written
(\S\ref{sec:audit:r04}): its qualifying sentence asserts negligible loss on COCO
CIDEr, a generation metric that assigns a graded score to a caption and has no
per-item correct/incorrect state, so V1 and V2, which are defined on $d_i \in
\{-1,0,+1\}$, have nothing to compute and there is no discordance rate to impute.
It is recorded as outside our registered calculation and carries no verdict. The
history of the overruled first-pass computation on a substituted benchmark is in
\S\ref{sec:audit:r04} and Appendix~\ref{app:audit:r04detail}, and that
computation is not claimed anywhere in this paper as an audit result.

The extension above does not retroactively score it, and it is worth being
precise about why, because the reason is not the mathematics. Equation
\eqref{eq:nreq-continuous} would need three inputs that do not exist here:
per-caption CIDEr scores for both models on the same images, from which
$\sigma_d$ could be measured; a margin declared in CIDEr points and defended as
a statement about caption quality; and a decision, taken before looking, about
which of those a claim of negligible loss is asserting.
The audit records that \AuditPerItemTaskMatched{} of the \AuditEligible{}
eligible sources release \emph{task-matched} per-item outputs
(\S\ref{sec:audit:v3}), so the first input is unavailable from the public
record, and the second and third were never registered.
The limit is the model we registered and the evidence the source released, not
paired analysis, which handles graded metrics without difficulty.
